% Λύση του 66ου προβλήματος της ενότητας «Άλυτα Προβλήματα».
\subsection*{Λύση Προβλήματος 66 --- Καμπύλη Gompertz και σημείο καμπής}

\begin{alist}
  \item Θέτουμε
        \[
          u(t)=e^{-0{,}4(t-10)}.
        \]
        Τότε
        \[
          G(t)=100e^{-u(t)}.
        \]
        Επίσης
        \[
          u'(t)=-0{,}4e^{-0{,}4(t-10)}=-0{,}4u(t).
        \]
        Άρα
        \[
          G'(t)
          =
          100e^{-u(t)}\cdot(-u'(t)).
        \]
        Επομένως
        \[
          G'(t)
          =
          100e^{-u(t)}\cdot0{,}4u(t)
          =
          40u(t)e^{-u(t)}.
        \]
        Δηλαδή
        \[
          {
          G'(t)=40e^{-0{,}4(t-10)}e^{-e^{-0{,}4(t-10)}}
          }.
        \]

  \item Από το προηγούμενο ερώτημα
        \[
          G'(t)=40u(t)e^{-u(t)}.
        \]
        Παραγωγίζουμε:
        \[
          G''(t)
          =
          40u'(t)e^{-u(t)}+40u(t)e^{-u(t)}(-u'(t)).
        \]
        Άρα
        \[
          G''(t)
          =
          40u'(t)e^{-u(t)}(1-u(t)).
        \]
        Επειδή \(u'(t)=-0{,}4u(t)\), παίρνουμε
        \[
          G''(t)
          =
          -16u(t)e^{-u(t)}(1-u(t))
          =
          16u(t)e^{-u(t)}(u(t)-1).
        \]
        Επομένως
        \[
          {
          G''(t)=16e^{-0{,}4(t-10)}e^{-e^{-0{,}4(t-10)}}
          \left(e^{-0{,}4(t-10)}-1\right)
          }.
        \]

  \item Για κάθε \(t\geq0\), ισχύει
        \[
          16e^{-0{,}4(t-10)}e^{-e^{-0{,}4(t-10)}}>0.
        \]
        Άρα το πρόσημο της \(G''(t)\) είναι το πρόσημο του
        \[
          e^{-0{,}4(t-10)}-1.
        \]

        Λύνουμε:
        \[
          e^{-0{,}4(t-10)}-1=0
          \quad\Longleftrightarrow\quad
          e^{-0{,}4(t-10)}=1.
        \]
        Άρα
        \[
          -0{,}4(t-10)=0
          \quad\Longleftrightarrow\quad
          t=10.
        \]

        Για \(0\leq t<10\), έχουμε
        \[
          e^{-0{,}4(t-10)}>1,
        \]
        οπότε \(G''(t)>0\). Για \(t>10\), έχουμε
        \[
          e^{-0{,}4(t-10)}<1,
        \]
        οπότε \(G''(t)<0\).

        Επομένως η κυρτότητα αλλάζει στο \(t=10\). Το σημείο καμπής
        είναι
        \[
          \left(10,G(10)\right).
        \]
        Υπολογίζουμε:
        \[
          G(10)=100e^{-e^0}=100e^{-1}=\frac{100}{e}.
        \]
        Άρα το σημείο καμπής είναι
        \[
          {
          \left(10,\frac{100}{e}\right)
          }
        \]
        δηλαδή προσεγγιστικά
        \[
          {(10,\ 36{,}8)}.
        \]

  \item Το σημείο καμπής είναι η στιγμή που ο ρυθμός αύξησης \(G'(t)\)
        γίνεται μέγιστος.

        Για \(0\leq t<10\), η \(G\) είναι κυρτή, άρα ο ρυθμός αύξησης
        αυξάνεται. Για \(t>10\), η \(G\) είναι κοίλη, άρα ο ρυθμός
        αύξησης μειώνεται. Βιολογικά, στο μοντέλο αυτό η ανάπτυξη
        επιταχύνεται μέχρι τη \(10\)η ημέρα και μετά συνεχίζει να
        αυξάνεται, αλλά με όλο και μικρότερο ρυθμό, επειδή πλησιάζει
        σε κορεσμό.
\end{alist}

\paragraph{Πηγές και έλεγχος ρεαλιστικότητας}
Η συνάρτηση \en{Gompertz}
\[
  f(t)=ae^{-be^{-ct}}
\]
είναι καθιερωμένη σιγμοειδής καμπύλη ανάπτυξης και χρησιμοποιείται σε
βιολογικά μοντέλα, όπως πληθυσμοί κυττάρων και μοντέλα ανάπτυξης όγκων
\cite{gompertzFunctionGrowth}. Το αποτέλεσμα είναι ρεαλιστικό ως
ποιοτική περιγραφή: η ανάπτυξη είναι αργή στην αρχή, μετά επιταχύνεται,
και στη συνέχεια επιβραδύνεται καθώς πλησιάζει σε ανώτατη τιμή. Στο
συγκεκριμένο κανονικοποιημένο μοντέλο το σημείο καμπής βρίσκεται στο
\[
  G=\frac{100}{e}\approx36{,}8\%
\]
της τελικής τιμής, χαρακτηριστικό της καμπύλης \en{Gompertz}.
